\chapter{Introducción}

\noindent\begin{minipage}[c]{0.76\textwidth}
Todo el material de este trabajo está publicado en un repositorio de acceso
libre. Incluye los guiones de preparación del conjunto de datos y de
entrenamiento del modelo, el firmware de las tres placas, los ficheros
paramétricos de la carcasa y las fuentes de este documento.

\smallskip
\noindent\url{https://github.com/xalsave/Sintetizador-de-Espacio-Latente}
\end{minipage}\hfill
\begin{minipage}[c]{0.20\textwidth}
\centering
\qrcode[height=2.4cm]{https://github.com/xalsave/Sintetizador-de-Espacio-Latente}
\end{minipage}

\bigskip

\section{Contexto y planteamiento del problema}
\hypertarget{RE1.1}{}\label{RE1.1}

La síntesis wavetable es una de las técnicas más extendidas en los instrumentos
musicales electrónicos actuales. Su principio es sencillo, en lugar de calcular
la forma de onda muestra a muestra mediante una expresión matemática, se
almacena un ciclo completo de la onda en una tabla y se recorre esa tabla a la
velocidad a la que se quiere reproducir la nota. El coste
computacional es bajo y el control sobre el timbre es directo, ya que basta con
sustituir el contenido de la tabla para obtener un sonido distinto.

El problema aparece al intentar pasar de un timbre a otro de forma continua. Un
banco de wavetables es un conjunto discreto, contiene las ondas que contiene y
nada más. Promediar dos de ellas sin ningún criterio previo no produce un timbre
intermedio convincente, sino cancelaciones entre armónicos que suenan a filtro
de peine, porque las dos ondas no comparten una referencia de fase común.
Tampoco existe una noción evidente de cercanía entre ondas. Dado un banco de
varios miles de formas de onda no hay una ordenación natural que sitúe juntas
las que suenan parecido.

Los autoencoders ofrecen una vía para resolver las dos cosas a la vez. Un
autoencoder es una red neuronal formada por dos mitades: un codificador, que
comprime cada onda a un número reducido de valores; y un decodificador, que
reconstruye la onda a partir de esos valores. El conjunto de valores intermedios
recibe el nombre de espacio latente. Si ese espacio se restringe a dos
dimensiones, cada onda de un conjunto queda situada en un punto de un plano, y las
ondas de sonido parecido tienden a agruparse. La variante variacional del
autoencoder añade una condición adicional sobre la organización de ese plano. Lo
obliga a ser continuo, sin huecos ni saltos bruscos. La consecuencia
práctica es que el decodificador puede generar una forma de onda con sentido para
cualquier punto del plano, incluidos los que no corresponden a ninguna onda del
conjunto original. Ahí es donde nacen los timbres intermedios que un banco
discreto no puede proporcionar.

Un plano continuo de dos dimensiones es, además, exactamente lo que un dedo
puede controlar sobre una pantalla táctil. Esa coincidencia entre la dimensión
del espacio aprendido y la del gesto de control es la que da sentido al
instrumento que se plantea en este trabajo.

\section{Motivación y oportunidad}
\hypertarget{RE1.1b}{}\label{RE1.1b}

La idea de generar wavetables mediante redes neuronales no es nueva. Existen
trabajos previos que aplican autoencoders a la generación e interpolación de
formas de onda de un ciclo, así como sistemas de síntesis neuronal en tiempo
real capaces de funcionar sobre hardware de propósito general. Todos ellos, sin
embargo, comparten un rasgo común: se materializan como aplicaciones de
escritorio o como complementos para estaciones de trabajo de audio digital. El
sonido se genera en un ordenador.

La oportunidad que aborda este trabajo es distinta. Se trata de llevar ese mismo
concepto a un instrumento físico autónomo basado en microcontroladores
de bajo coste que generen el audio en su propio hardware y respondan a un gesto
táctil y a un teclado \gls{midi} sin necesidad de un ordenador conectado. La
aportación no está, por tanto, en la interpolación neuronal en sí, sino en su
integración embebida completa en tiempo real. En resolver el reparto de tareas
entre procesadores, la comunicación entre ellos, la latencia, la estabilidad del
audio y la construcción del prototipo.

Esa integración impone una restricción que condiciona todo el diseño y que
conviene enunciar desde el principio: la red neuronal no puede ejecutarse a
frecuencia de audio. Generar 48.000 muestras por segundo mediante inferencia
neuronal queda fuera del alcance de un microcontrolador de estas
características. Ahora bien, la red no necesita trabajar a esa velocidad. Solo
tiene que producir una forma de onda nueva cuando el usuario mueve el dedo, es
decir, unas pocas decenas de veces por segundo. La reproducción de esa onda a
48~kHz es una operación aritmética elemental. Separar la generación de la tabla,
lenta y esporádica, de su reproducción, rápida y continua, es lo que hace viable
el sistema completo.

Desde el punto de vista académico, el trabajo integra contenidos de
procesamiento digital de señal, aprendizaje automático, sistemas embebidos,
protocolos de comunicación entre dispositivos y diseño electrónico. Es una gran
oportunidad para aplicar transversalmente todos los conocimientos y competencias
adquiridos en el grado.

\section{Objetivos}
\hypertarget{RE1.3}{}\label{RE1.3}

El objetivo general es diseñar, construir y validar un sintetizador digital
embebido y autónomo que permita explorar de forma continua un espacio tímbrico
aprendido por un autoencoder, empleando una pantalla táctil como superficie de
control y una entrada \gls{midi} estándar para la ejecución musical.

De este objetivo general se derivan los siguientes objetivos específicos:

\begin{itemize}
    \item Construir un conjunto de datos homogéneo a partir de un conjunto de datos público
          de formas de onda de un ciclo, resolviendo el remuestreo, la
          normalización de amplitud y la alineación de fase necesarias para que
          la interpolación posterior no introduzca artefactos.
    \item Entrenar un autoencoder variacional que comprima cada forma de onda a
          una coordenada bidimensional y verificar que el espacio latente
          resultante es continuo y está bien aprovechado.
    \item Precalcular, a partir del decodificador entrenado, una rejilla de
          formas de onda que pueda embarcarse en el sistema y exportarla en un
          formato adecuado para un microcontrolador.
    \item Implementar un motor de síntesis wavetable en tiempo real capaz de
          sustituir la tabla en reproducción sin discontinuidades audibles.
    \item Resolver la comunicación entre los tres microcontroladores que
          componen el sistema, definiendo los protocolos y verificando que la
          información llega íntegra.
    \item Incorporar el control de nota y dinámica mediante \gls{midi}, así como un
          panel de control analógico para los parámetros de envolvente y filtro.
    \item Validar cada etapa de forma objetiva, contrastando el comportamiento
          del sistema embebido con una implementación de referencia en el
          ordenador siempre que sea posible.
    \item Integrar el conjunto en un prototipo físico alimentado desde una
          fuente única, con su propia electrónica de soporte.
\end{itemize}

\section{Alcance y condicionantes}
\hypertarget{RE1.2}{}\label{RE1.2}

El sistema se construye sobre tres microcontroladores con funciones separadas.
Una pantalla táctil con controlador ESP32 integrado se encarga de la interfaz de
usuario. Un ESP32-S3 aloja la rejilla de formas de onda y calcula la
interpolación. Una placa Daisy Seed, basada en un procesador ARM Cortex-M7,
genera el audio gracias a su \gls{dac} de 24 bits. El teclado que proporciona las notas se conecta por MIDI DIN y
se considera un elemento externo e intercambiable, de forma que el entregable
del trabajo es un módulo de sonido y no un instrumento con teclado integrado.

El alcance queda acotado por varios condicionantes que conviene dejar
explícitos:

\begin{itemize}
    \item \textbf{La red neuronal no se ejecuta en el instrumento.} El
          decodificador se emplea fuera de línea, en el ordenador, para generar
          por adelantado la rejilla de formas de onda que se embarca en el
          sistema. Durante la ejecución solo hay consulta de tabla e
          interpolación. Se trata de una estrategia de despliegue, no de una
          renuncia: es la consecuencia directa de la restricción enunciada en el
          apartado anterior.
    \item \textbf{El espacio latente es bidimensional.} Esta elección viene
          impuesta por el control táctil, pero supone un cuello de botella
          considerable para representar varios miles de formas de onda
          distintas. Su efecto sobre la fidelidad de la reconstrucción se mide y
          se discute en el capítulo sobre el modelo neuronal.
    \item \textbf{El instrumento es monofónico.} La polifonía no forma parte de
          los objetivos y se plantea únicamente como línea de trabajo futuro.
    \item \textbf{El hardware es de bajo coste y de disponibilidad comercial.}
          No se ha diseñado ni fabricado ningún circuito integrado específico, y
          el montaje final se resuelve sobre placa de tiras de cobre en lugar de
          sobre un circuito impreso fabricado, por restricciones de plazo.
    \item \textbf{La amplificación y el altavoz integrados quedan fuera del
          alcance.} La salida del instrumento es una señal de línea por conector
          de 3,5~mm. 
\end{itemize}

\section{Estructura de la memoria}

El presente documento se organiza en nueve capítulos.

El \textbf{capítulo 1} sitúa el problema, expone la motivación del trabajo, fija
los objetivos y delimita el alcance.

El \textbf{capítulo 2} recoge los fundamentos teóricos necesarios para seguir el
resto del documento y revisa los trabajos previos relacionados con
sistemas de síntesis neuronal.

El \textbf{capítulo 3} describe el diseño del sistema completo: el reparto de
funciones entre los tres microcontroladores, los protocolos de comunicación
entre ellos y la separación entre la fase de preparación fuera de línea y la de
ejecución en tiempo real.

El \textbf{capítulo 4} aborda la parte de aprendizaje automático, desde el
preprocesado del conjunto de formas de onda hasta el entrenamiento del
autoencoder y la generación de la rejilla que se embarca en el sistema.

Los \textbf{capítulos 5 y 6} describen el desarrollo del firmware. El primero se
ocupa de la cadena de entrada, formada por la pantalla táctil y el ESP32-S3, y
el segundo del motor de audio implementado sobre la Daisy Seed.

El \textbf{capítulo 7} trata la parte física del prototipo: la alimentación, el
panel de control analógico y el montaje definitivo.

El \textbf{capítulo 8} reúne las pruebas realizadas y los resultados obtenidos,
tanto las medidas eléctricas como las verificaciones numéricas de cada etapa de
la cadena de procesado.

El \textbf{capítulo 9} recoge las conclusiones y plantea las líneas de trabajo
futuro.
